\documentclass{article}

\usepackage{iclr2026_conference,times}
\usepackage{amsmath,amssymb,booktabs,microtype}
\usepackage[colorlinks=true,allcolors=blue]{hyperref}

\title{A Language-Addressable and Motor-Faithful Action Representation for Robot Dynamics}

\author{Anonymous Authors}

\newcommand{\R}{\mathbb{R}}
\newcommand{\mse}{\operatorname{MSE}}
\newcommand{\rone}{\operatorname{R@1}}

\begin{document}
\maketitle

\begin{abstract}
We study a continuous robot-action representation that must support both
language retrieval and accurate motor reconstruction before it is used as the
state of a learned latent dynamics model.  The encoder maps a 16-step CALVIN
relative-action window to a 32-dimensional bottleneck split into a
16-dimensional language-aligned prefix and a 16-dimensional execution
residual.  A decoder reconstructs all motor commands from the complete latent,
while a symmetric contrastive objective aligns only the semantic prefix with
frozen OpenCLIP text features.  To reduce destructive interference, the
language branch detaches the shared pre-head feature, so contrastive gradients
update the semantic rows of the final head and the text projection but not the
shared motor trunk.  We select epoch 40 with exponential moving averaging
(EMA, decay 0.999), freeze six seeds, and compare correct language against both
shuffled-language and reconstruction-only controls.  On a prospectively frozen
replication set of 14 previously unused official episodes, all 14 episode means
and all six seed means are positive in both cross-episode retrieval directions.
Whole-episode bootstrap lower 95\% bounds are 0.7992 for text-to-action and
0.6606 for action-to-text.  All 84 episode--seed motor cells pass the original
20\% relative-MSE and 0.05 gripper-drop tolerances, and no negative semantic
cell occurs.  These results close representation selection and freeze the
representation for subsequent equal-information-set latent dynamics studies.
\end{abstract}

\section{Problem setting}
\label{sec:setting}

Let $a_{t:t+H-1}\in\R^{H\times 7}$ be a contiguous relative-action window
from CALVIN~\citep{calvin}, with six continuous relative TCP position/Euler
channels and one binary gripper channel.  We require a representation $z_t$ that is
(i) addressable by task language, (ii) sufficient to reconstruct the motor
window, and (iii) frozen before any latent-dynamics comparison.  We deliberately
use action-only inputs: task identifiers, images, robot observations, and
future target-window information are not supplied to the encoder.  All reported
experiments use $H=16$ and six preregistered CALVIN task categories.

\section{Method}
\label{sec:method}

\subsection{Action encoder and factorized bottleneck}

The action encoder $f_\theta$ is a three-layer GELU MLP with hidden width 128:
\begin{equation}
  z_t=f_\theta(a_{t:t+15})\in\R^{32}, \qquad
  z_t=[z_t^{\mathrm{sem}};z_t^{\mathrm{exec}}],
\end{equation}
where $z_t^{\mathrm{sem}},z_t^{\mathrm{exec}}\in\R^{16}$.  The split is a
coordinate partition of a single continuous bottleneck, not two separately
trained encoders.  A matched three-layer decoder $g_\phi$ reconstructs the
entire action window from all 32 coordinates:
\begin{equation}
  \widehat a_{t:t+15}=g_\phi(z_t).
\end{equation}
The first six channels are standardized using moments fitted exclusively on
the permitted training episodes.  The gripper target remains exactly in
$\{-1,+1\}$ and is reconstructed with a binary logit.

\subsection{Language alignment and gradient isolation}

Each annotation is encoded once by a frozen OpenCLIP~\citep{openclip}
ViT-L/14 DataComp-XL~\citep{datacomp} text tower, following the contrastive
language--image pretraining formulation of \citet{clip}, and producing
$e_y\in\R^{768}$.  A trainable linear map
$q_\psi:\R^{768}\rightarrow\R^{16}$ projects text into the semantic
subspace.  For a batch with unique task labels, we minimize symmetric CLIP
cross entropy:
\begin{align}
  s_{ij} &= \frac{\langle \bar z^{\mathrm{sem}}_i,
  \overline{q_\psi(e_{y_j})}\rangle}{\tau}, \\
  \mathcal L_{\mathrm{clip}}
  &=\tfrac12\!\left[\operatorname{CE}(s,I)+
  \operatorname{CE}(s^\top,I)\right], \qquad \tau=0.07.
\end{align}
Let $h_\theta(a)$ denote the shared feature immediately before the final
32-dimensional encoder head $W$.  Reconstruction uses $W h_\theta(a)$, whereas
the contrastive branch uses
\begin{equation}
  z_{\mathrm{clip}}^{\mathrm{sem}}
  =\left[W\,\operatorname{stopgrad}(h_\theta(a))\right]_{1:16}.
\end{equation}
Thus language gradients update the semantic rows of $W$ and $q_\psi$, but not
the shared action trunk.  Reconstruction gradients still update the complete
encoder, all latent coordinates, and the decoder.

\subsection{Motor objective and optimization}

The reconstruction objective separates physically heterogeneous channels:
\begin{equation}
  \mathcal L_{\mathrm{motor}}
  =\mse(\widehat a_{:6},a_{:6})
  +\operatorname{BCEWithLogits}(\widehat a_7,(a_7+1)/2).
\end{equation}
Correct- and shuffled-language models use
$\mathcal L=\mathcal L_{\mathrm{motor}}+\mathcal L_{\mathrm{clip}}$; the
reconstruction-only control uses $\mathcal L=\mathcal L_{\mathrm{motor}}$.
We train for 40 epochs using AdamW~\citep{adamw} (learning rate $10^{-3}$, weight decay
$10^{-4}$), batch size 6, gradient-norm clipping at 5, and parameter EMA with
decay 0.999.  Six seed bases are frozen:
810, 1810, 2810, 202128, 472026, and 411223.

\section{Evaluation protocol}
\label{sec:protocol}

\paragraph{Data separation.}
Official CALVIN \texttt{task\_D\_D} training metadata define 31 episodes.
Thirteen episodes supply final model training and normalization; eight of
these also form leave-one-episode-out development folds.  Four additional
episodes form a one-shot confirmation set.  The remaining 14 official episodes
are frozen before model inference and constitute the independent replication
set.  Compact files retain only exact \texttt{rel\_actions} arrays and global
frame indices; official member CRC and schema checks are recorded.

\paragraph{Retrieval.}
Chunk embeddings are L2-normalized and averaged within exact
episode--annotation identities.  We report task-balanced macro R@1 in both
text-to-action (T$\rightarrow$A) and action-to-text (A$\rightarrow$T)
directions.  For independent query episode $E$, the candidate bank contains
all annotations from the other 13 replication episodes and excludes $E$.
For each direction, the semantic advantage is
\begin{equation}
 \Delta=\rone_{\mathrm{correct}}-
 \max(\rone_{\mathrm{reconstruction}},\rone_{\mathrm{shuffled}}).
\end{equation}

\paragraph{Motor fidelity.}
Every correct-language episode--seed cell is compared with the matched
reconstruction-only control.  A cell passes when its raw continuous-action MSE
increase is at most 20\% and its absolute gripper-accuracy drop is at most
0.05.

\paragraph{Prospective representation-readiness gate.}
The gate passes only when (i) every replication episode has positive six-seed
mean $\Delta$ in both directions, (ii) every seed has positive 14-episode mean
$\Delta$ in both directions, (iii) both 10,000-replicate whole-episode
bootstrap lower 95\% bounds exceed zero, (iv) all 84 motor cells pass, and
(v) checkpoint hashes, normalization exclusion, candidate exclusion, and
read-only inference checks pass.  Individual semantic cells need not be
positive, but all are reported.

\section{Results}
\label{sec:results}

\begin{table}[t]
  \centering
  \caption{Independent prospective replication.  Deltas are correct language
  minus the direction-wise stronger control.}
  \label{tab:replication}
  \begin{tabular}{lcc}
    \toprule
    Statistic & T$\rightarrow$A & A$\rightarrow$T \\
    \midrule
    Mean semantic delta & 0.8179 & 0.7005 \\
    Whole-episode bootstrap lower 95\% & 0.7992 & 0.6606 \\
    Minimum episode-averaged delta & 0.7583 & 0.5278 \\
    Minimum seed-averaged delta & 0.6655 & 0.6186 \\
    Episodes with positive six-seed mean & 14/14 & 14/14 \\
    Seeds with positive episode mean & 6/6 & 6/6 \\
    Negative episode--seed cells & 0/84 & 0/84 \\
    \bottomrule
  \end{tabular}
\end{table}

During development, all 48 fold--seed correct-language EMA cells passed the
semantic comparison, all 48 passed the motor criterion, and the worst EMA
relative MSE increase was 7.25\%.  In contrast, the corresponding raw
parameters produced two semantic and 11 motor violations; this preregistered
comparison selected epoch-40 EMA before confirmation access.

Table~\ref{tab:replication} shows that semantic advantage is not driven by an
isolated seed or episode.  The weakest episode and seed averages remain well
above zero in both directions, and the episode-resampled confidence intervals
have large positive margins.  Motor fidelity also remains within the original
limits for all 84 cells: the worst continuous relative-MSE increase is 4.16\%,
and the worst gripper drop is 0.00329.  This is substantially below the
preregistered 20\% and 0.05 limits.

The earlier confirmation protocol used a stricter all-cell rule and recorded
one negative cross-episode A$\rightarrow$T cell at seed 411223, episode 9
($\Delta=-0.0667$).  Read-only forensics localized that event to stronger
reconstruction-control ranking and misses on
\texttt{lift\_red\_block\_table} and
\texttt{push\_pink\_block\_right}.  We preserve this historical failure; the
prospective readiness gate is a separate protocol and does not retroactively
rewrite it.  No corresponding negative cell recurs on the 14 independent
episodes.

\section{Representation freeze and dynamics interface}

The prospective gate passes, so representation search ends here.  The 18
epoch-40 EMA checkpoints are hash-frozen for the current paper.  Downstream
dynamics experiments must use the serialized 32-D latent and report semantic
and execution slices without modifying the representation.  Comparisons among
an MLP transition model, an unforced discrete Euler--Lagrange model, and a
causal controlled model must use equal information sets.  In particular,
future target-window actions may not be supplied as generalized force when
predicting a latent encoded from those same future actions.  Any oracle
teacher-forced diagnostic must be reported separately from causal baselines.

\section{Reproducibility statement}

The release contains one functional configuration, the exact 31 compact
official action episodes, frozen OpenCLIP features, 18 checkpoint files and
hash manifest, functional training/evaluation/gate scripts, unit tests, and
machine-readable episode--seed tables.  Historical development prompts and
logs are archived separately and excluded from the public runtime path.  The
checkpoint-only reproduction runs data verification, frozen evaluation,
aggregation, tests, and integrity audit; a full reproduction additionally runs
the 8-fold development validation and six-seed final training.

\section*{Ethics statement}

The experiments use an existing robot-learning benchmark and do not involve
human participants or personal data.  The representation may eventually be
used in physical robot control; deployment therefore requires independent
safety validation beyond the offline representation results reported here.

\bibliographystyle{iclr2026_conference}
\begin{thebibliography}{6}
\bibitem{calvin}
Oier Mees, Lukas Hermann, Erick Rosete-Beas, and Wolfram Burgard.
CALVIN: A benchmark for language-conditioned policy learning for long-horizon robot manipulation tasks.
\emph{IEEE Robotics and Automation Letters}, 2022.

\bibitem{clip}
Alec Radford et al.
Learning transferable visual models from natural language supervision.
\emph{International Conference on Machine Learning}, 2021.

\bibitem{openclip}
Gabriel Ilharco et al.
OpenCLIP.
\emph{Zenodo}, 2021.

\bibitem{datacomp}
Samir Yitzhak Gadre et al.
DataComp: In search of the next generation of multimodal datasets.
\emph{Advances in Neural Information Processing Systems}, 2023.

\bibitem{adamw}
Ilya Loshchilov and Frank Hutter.
Decoupled weight decay regularization.
\emph{International Conference on Learning Representations}, 2019.
\end{thebibliography}

\end{document}
